% GENERATED FILE. Edit canonical lessons, then run npm run generate:books.
% canonical-source-hash: fnv1a64:86bb4d29098c8600
% canonical-lessons: TA-C123-kaykari, TA-C123-valaippalam, TA-C123-mampalam, TA-C123-tenkay, TA-C123-venkayam

\chapter{Vegetables, Banana, Mango, Coconut, Onion}
\label{ch:ta-vegetables-banana-mango-coconut-onion}
\hlchaptermodality{123}

\begin{chapteropening}
\textbf{By the end of this chapter:} \emph{I can say vegetables, a banana, a mango, a coconut and an onion.}

Complete the last lesson of chapter 123: i can say vegetables, a banana, a mango, a coconut and an onion.
\end{chapteropening}

\section[\texorpdfstring{\ta{காய்கறி} (kāykaṟi)}{காய்கறி (kāykaṟi)}]{\ta{காய்கறி} (kāykaṟi) — vegetables}

\label{lesson:TA-C123-kaykari}



\begin{quote}
Before the new one: say the Tamil for a fish, then the Tamil for meat.
\end{quote}

\subsection*{You'll want to know: \ta{காய்கறி}}
\textbf{\ta{காய்கறி}} — \emph{kāykaṟi} — \textquotedblleft{}vegetables\textquotedblright{}.

\textbf{In the family, and next door}

\noindent
\begin{tabularx}{\linewidth}{@{}>{\raggedright\arraybackslash}X>{\raggedright\arraybackslash}X@{}}
\toprule
\textbf{} & \textbf{word} \\
\midrule
Kannada & \ka{ತರಕಾರಿ} (\emph{tarakāri}) \\
Telugu & \te{కూరగాయలు} (\emph{kūragāyalu}) \\
Malayalam & \ml{പച്ചക്കറി} (\emph{paccakkaṟi}) \\
English & vegetables \\
\bottomrule
\end{tabularx}

\begin{grammarlens}[title={What you've built}]
One more word for this chapter.
\end{grammarlens}

\subsection*{Your turn}
\begin{itemize}
\raggedright
  \item \emph{Say it:} \emph{kāykaṟi}
  \item \emph{Say it:} \emph{kāykaṟi}, once more
  \item \emph{Say it:} say \emph{kāykaṟi}
  \item \emph{Recall:} say the Tamil for a fish, then the Tamil for meat, then say \emph{kāykaṟi} again
\end{itemize}

\begin{tcolorbox}[breakable,colback=teal!4,colframe=teal!35!black,title={Before you move on}]
What does \ta{காய்கறி} mean? (\textquotedblleft{}Vegetables\textquotedblright{}.) Say it once more.
\end{tcolorbox}
\section[\texorpdfstring{\ta{வாழைப்பழம்} (vāḻaippaḻam)}{வாழைப்பழம் (vāḻaippaḻam)}]{\ta{வாழைப்பழம்} (vāḻaippaḻam) — a banana}

\label{lesson:TA-C123-valaippalam}



\begin{quote}
Before the new one: say the Tamil for meat, then the Tamil for vegetables.
\end{quote}

\subsection*{You'll want to know: \ta{வாழைப்பழம்}}
\textbf{\ta{வாழைப்பழம்}} — \emph{vāḻaippaḻam} — \textquotedblleft{}a banana\textquotedblright{}.

\textbf{In the family, and next door}

\noindent
\begin{tabularx}{\linewidth}{@{}>{\raggedright\arraybackslash}X>{\raggedright\arraybackslash}X>{\raggedright\arraybackslash}X@{}}
\toprule
\textbf{} & \textbf{word} & \textbf{same root} \\
\midrule
Kannada & \ka{ಬಾಳೆಹಣ್ಣು} (\emph{bāḷehaṇṇu}) &  \\
Telugu & \te{అరటిపండు} (\emph{araṭipaṇḍu}) &  \\
Malayalam & \ml{വാഴപ്പഴം} (\emph{vāḻappaḻaṁ}) & yes \\
Hindi (neighbour) & \dv{केला} (\emph{kelā}) &  \\
English & a banana &  \\
\bottomrule
\end{tabularx}

\begin{grammarlens}[title={What you've built}]
One more word for this chapter.
\end{grammarlens}

\subsection*{Your turn}
\begin{itemize}
\raggedright
  \item \emph{Say it:} \emph{vāḻaippaḻam}
  \item \emph{Say it:} \emph{vāḻaippaḻam}, once more
  \item \emph{Say it:} say \emph{vāḻaippaḻam}
  \item \emph{Recall:} say the Tamil for meat, then the Tamil for vegetables, then say \emph{vāḻaippaḻam} again
\end{itemize}

\begin{tcolorbox}[breakable,colback=teal!4,colframe=teal!35!black,title={Before you move on}]
What does \ta{வாழைப்பழம்} mean? (\textquotedblleft{}A banana\textquotedblright{}.) Say it once more.
\end{tcolorbox}
\section[\texorpdfstring{\ta{மாம்பழம்} (māmpaḻam)}{மாம்பழம் (māmpaḻam)}]{\ta{மாம்பழம்} (māmpaḻam) — a mango}

\label{lesson:TA-C123-mampalam}



\begin{quote}
Before the new one: say the Tamil for vegetables, then the Tamil for a banana.

Say \textbf{\ta{எனக்குத்} \ta{தமிழ்} \ta{பிடிக்கும்}}, then the plain verb for wanting, \textbf{\ta{விரும்பு}}.
\end{quote}

\subsection*{You'll want to know: \ta{மாம்பழம்}}
\textbf{\ta{மாம்பழம்}} — \emph{māmpaḻam} — \textquotedblleft{}a mango\textquotedblright{}.

\textbf{In the family, and next door}

\noindent
\begin{tabularx}{\linewidth}{@{}>{\raggedright\arraybackslash}X>{\raggedright\arraybackslash}X>{\raggedright\arraybackslash}X@{}}
\toprule
\textbf{} & \textbf{word} & \textbf{same root} \\
\midrule
Kannada & \ka{ಮಾವಿನಹಣ್ಣು} (\emph{māvinahaṇṇu}) &  \\
Telugu & \te{మామిడిపండు} (\emph{māmiḍipaṇḍu}) & yes \\
Hindi (neighbour) & \dv{आम} (\emph{ām}) &  \\
English & a mango &  \\
\bottomrule
\end{tabularx}

\begin{grammarlens}[title={What you've built}]
One more word for this chapter.
\end{grammarlens}

\subsection*{Your turn}
\begin{itemize}
\raggedright
  \item \emph{Say it:} \emph{māmpaḻam}
  \item \emph{Say it:} \emph{māmpaḻam}, once more
  \item \emph{Say it:} say \emph{māmpaḻam}
  \item \emph{Recall:} say the Tamil for vegetables, then the Tamil for a banana, then say \emph{māmpaḻam} again
\end{itemize}

\begin{tcolorbox}[breakable,colback=teal!4,colframe=teal!35!black,title={Before you move on}]
What does \ta{மாம்பழம்} mean? (\textquotedblleft{}A mango\textquotedblright{}.) Say it once more.
\end{tcolorbox}
\section[\texorpdfstring{\ta{தேங்காய்} (tēṅkāy)}{தேங்காய் (tēṅkāy)}]{\ta{தேங்காய்} (tēṅkāy) — a coconut}

\label{lesson:TA-C123-tenkay}



\begin{quote}
Before the new one: say the Tamil for a banana, then the Tamil for a mango.
\end{quote}

\subsection*{You'll want to know: \ta{தேங்காய்}}
\textbf{\ta{தேங்காய்}} — \emph{tēṅkāy} — \textquotedblleft{}a coconut\textquotedblright{}.

\textbf{In the family, and next door}

\noindent
\begin{tabularx}{\linewidth}{@{}>{\raggedright\arraybackslash}X>{\raggedright\arraybackslash}X>{\raggedright\arraybackslash}X@{}}
\toprule
\textbf{} & \textbf{word} & \textbf{same root} \\
\midrule
Kannada & \ka{ತೆಂಗು} (\emph{teṅgu}) & yes \\
Telugu & \te{కొబ్బరి} (\emph{kobbari}) &  \\
Malayalam & \ml{തേങ്ങ} (\emph{tēṅṅa}) & yes \\
Hindi (neighbour) & \dv{नारियल} (\emph{nāriyal}) &  \\
English & a coconut &  \\
\bottomrule
\end{tabularx}

\begin{grammarlens}[title={What you've built}]
One more word for this chapter.
\end{grammarlens}

\subsection*{Your turn}
\begin{itemize}
\raggedright
  \item \emph{Say it:} \emph{tēṅkāy}
  \item \emph{Say it:} \emph{tēṅkāy}, once more
  \item \emph{Say it:} say \emph{tēṅkāy}
  \item \emph{Recall:} say the Tamil for a banana, then the Tamil for a mango, then say \emph{tēṅkāy} again
\end{itemize}

\begin{tcolorbox}[breakable,colback=teal!4,colframe=teal!35!black,title={Before you move on}]
What does \ta{தேங்காய்} mean? (\textquotedblleft{}A coconut\textquotedblright{}.) Say it once more.
\end{tcolorbox}
\section[\texorpdfstring{\ta{வெங்காயம்} (veṅkāyam)}{வெங்காயம் (veṅkāyam)}]{\ta{வெங்காயம்} (veṅkāyam) — an onion}

\label{lesson:TA-C123-venkayam}



\begin{quote}
Before the new one: say the Tamil for a mango, then the Tamil for a coconut.
\end{quote}

\subsection*{You'll want to know: \ta{வெங்காயம்}}
\textbf{\ta{வெங்காயம்}} — \emph{veṅkāyam} — \textquotedblleft{}an onion\textquotedblright{}.

\textbf{In the family, and next door}

\noindent
\begin{tabularx}{\linewidth}{@{}>{\raggedright\arraybackslash}X>{\raggedright\arraybackslash}X@{}}
\toprule
\textbf{} & \textbf{word} \\
\midrule
Kannada & \ka{ಈರುಳ್ಳಿ} (\emph{īruḷḷi}) \\
Telugu & \te{ఉల్లిపాయ} (\emph{ullipāya}) \\
Malayalam & \ml{ഉള്ളി} (\emph{uḷḷi}) \\
Hindi (neighbour) & \dv{प्याज़} (\emph{pyāz}) \\
English & an onion \\
\bottomrule
\end{tabularx}

\begin{grammarlens}[title={What you've built}]
One more word for this chapter.
\end{grammarlens}

\subsection*{Your turn}
\begin{itemize}
\raggedright
  \item \emph{Say it:} \emph{veṅkāyam}
  \item \emph{Say it:} \emph{veṅkāyam}, once more
  \item \emph{Say it:} say \emph{veṅkāyam}
  \item \emph{Recall:} say the Tamil for a mango, then the Tamil for a coconut, then say \emph{veṅkāyam} again
\end{itemize}

\begin{tcolorbox}[breakable,colback=teal!4,colframe=teal!35!black,title={Before you move on}]
What does \ta{வெங்காயம்} mean? (\textquotedblleft{}An onion\textquotedblright{}.) Say it once more.
\end{tcolorbox}
